% ============================================================
%  TALLER DE PYTHON — Documento Principal
%  Grupo: [Nombre del Grupo]
%  Integrantes:
%    - Integrante 1 (secciones 1-2)
%    - Integrante 2 (secciones 3-4)
%    - Integrante 3 (secciones 5 + portada/conclusiones)
%
%  Para compilar: pdflatex main.tex  (dos veces para índice)
% ============================================================

\documentclass[12pt, a4paper]{article}

% ── Paquetes ────────────────────────────────────────────────
\usepackage[utf8]{inputenc}
\usepackage[T1]{fontenc}
\usepackage[spanish]{babel}
\usepackage{geometry}
\usepackage{graphicx}
\usepackage{hyperref}
\usepackage{listings}       % bloques de código
\usepackage{xcolor}         % colores para el código
\usepackage{amsmath}
\usepackage{fancyhdr}       % encabezados y pies de página
\usepackage{titlesec}       % estilos de secciones
\usepackage{booktabs}       % tablas profesionales
\usepackage{enumitem}       % listas personalizadas
\usepackage{tcolorbox}      % cajas de respuesta
\usepackage{setspace}

% ── Geometría de página ─────────────────────────────────────
\geometry{
    top    = 2.5cm,
    bottom = 2.5cm,
    left   = 3cm,
    right  = 2.5cm
}

% ── Estilo de código Python ─────────────────────────────────
\definecolor{codegreen}{rgb}{0.1,0.6,0.1}
\definecolor{codegray}{rgb}{0.5,0.5,0.5}
\definecolor{codepurple}{rgb}{0.58,0,0.82}
\definecolor{backcolour}{rgb}{0.97,0.97,0.97}

\lstdefinestyle{pythonstyle}{
    language         = Python,
    backgroundcolor  = \color{backcolour},
    commentstyle     = \color{codegreen},
    keywordstyle     = \color{blue}\bfseries,
    numberstyle      = \tiny\color{codegray},
    stringstyle      = \color{codepurple},
    basicstyle       = \ttfamily\small,
    breakatwhitespace= false,
    breaklines       = true,
    keepspaces       = true,
    numbers          = left,
    numbersep        = 8pt,
    showspaces       = false,
    showstringspaces = false,
    frame            = single,
    rulecolor        = \color{gray!40},
    tabsize          = 4,
    captionpos       = b,
}
\lstset{style=pythonstyle}

% ── Caja para respuestas ────────────────────────────────────
\tcbuselibrary{skins, breakable}
\newtcolorbox{respuesta}[1][]{
    colback  = blue!3!white,
    colframe = blue!40!black,
    fonttitle= \bfseries,
    title    = Respuesta,
    breakable,
    #1
}

% ── Encabezado y pie de página ──────────────────────────────
\pagestyle{fancy}
\fancyhf{}
\rhead{\textit{Taller de Python}}
\lhead{\textit{[Nombre del Grupo]}}
\cfoot{\thepage}
\renewcommand{\headrulewidth}{0.4pt}

% ── Metadatos del PDF ───────────────────────────────────────
\hypersetup{
    colorlinks = true,
    linkcolor  = blue!70!black,
    urlcolor   = blue,
    pdftitle   = {Taller de Python},
    pdfauthor  = {Grupo X},
}

% ============================================================
\begin{document}

% ── Portada ─────────────────────────────────────────────────
% ============================================================
%  PORTADA — Editar: Integrante 3
% ============================================================

\begin{titlepage}
    \centering
    \vspace*{2cm}

    {\LARGE \textbf{Universidad / Institución}}\\[0.5cm]
    {\large Facultad / Departamento}\\[0.3cm]
    {\large Asignatura: Programación en Python}

    \vspace{2cm}
    \rule{\linewidth}{0.5mm}\\[0.4cm]
    {\huge \bfseries Taller de Python}\\[0.2cm]
    \rule{\linewidth}{0.5mm}

    \vspace{2cm}

    \begin{tabular}{ll}
        \textbf{Integrante 1:} & Nombre Completo \\[0.3cm]
        \textbf{Integrante 2:} & Nombre Completo \\[0.3cm]
        \textbf{Integrante 3:} & Nombre Completo \\
    \end{tabular}

    \vspace{2cm}

    \begin{tabular}{ll}
        \textbf{Docente:}  & Nombre del Profesor \\[0.3cm]
        \textbf{Grupo:}    & Número de grupo     \\[0.3cm]
        \textbf{Fecha:}    & \today              \\
    \end{tabular}

    \vfill
    {\large Ciudad, País}
\end{titlepage}
\newpage


% ── Tabla de contenidos ─────────────────────────────────────
\tableofcontents
\newpage

% ── Secciones del taller ────────────────────────────────────
%  Cada integrante edita sus propios archivos:
%  Integrante 1 → seccion1.tex, seccion2.tex
%  Integrante 2 → seccion3.tex, seccion4.tex
%  Integrante 3 → seccion5.tex, conclusiones.tex

% ============================================================
%  SECCIÓN 1 — Responsable: Integrante 1
%  Título: [Título del ejercicio según el taller]
% ============================================================

\section{[Título Sección 1]}

% ── Enunciado ───────────────────────────────────────────────
\subsection*{Enunciado}

\begin{tcolorbox}[colback=gray!8, colframe=gray!50, title=Ejercicio 1]
    % Copiar aquí el enunciado del ejercicio tal como aparece en el taller.
    Escriba aquí el enunciado del ejercicio 1.
\end{tcolorbox}

% ── Análisis / Planteamiento ────────────────────────────────
\subsection{Análisis del problema}

Descripción breve del enfoque o estrategia utilizada para resolver el ejercicio.

% ── Código ──────────────────────────────────────────────────
\subsection{Código}

\begin{lstlisting}[caption={Solución ejercicio 1}, label={lst:ej1}]
# Escribir aquí el código Python de la solución

def solucion():
    pass
\end{lstlisting}

% ── Respuesta / Resultado ───────────────────────────────────
\subsection{Resultados}

\begin{respuesta}
    Describir aquí los resultados obtenidos al ejecutar el programa,
    incluyendo salidas por consola o capturas de pantalla si aplica.
\end{respuesta}

% ── (Opcional) Captura de pantalla ─────────────────────────
% \begin{figure}[h!]
%     \centering
%     \includegraphics[width=0.75\linewidth]{imagenes/resultado_ej1.png}
%     \caption{Salida del ejercicio 1}
%     \label{fig:ej1}
% \end{figure}

\newpage

% ============================================================
%  SECCIÓN 2 — Responsable: Integrante 1
2
%  Título: [Título del ejercicio según el taller]
% ============================================================

\section{[Título Sección 2]}

% ── Enunciado ───────────────────────────────────────────────
\subsection*{Enunciado}

\begin{tcolorbox}[colback=gray!8, colframe=gray!50, title=Ejercicio 2]
    Escriba aquí el enunciado del ejercicio 2.
\end{tcolorbox}

% ── Análisis / Planteamiento ────────────────────────────────
\subsection{Análisis del problema}

Descripción breve del enfoque o estrategia utilizada para resolver el ejercicio.

% ── Código ──────────────────────────────────────────────────
\subsection{Código}

\begin{lstlisting}[caption={Solución ejercicio 2}, label={lst:ej2}]
# Escribir aquí el código Python de la solución

def solucion():
    pass
\end{lstlisting}

% ── Respuesta / Resultado ───────────────────────────────────
\subsection{Resultados}

\begin{respuesta}
    Describir aquí los resultados obtenidos al ejecutar el programa.
\end{respuesta}

% ── (Opcional) Captura de pantalla ─────────────────────────
% \begin{figure}[h!]
%     \centering
%     \includegraphics[width=0.75\linewidth]{imagenes/resultado_ej2.png}
%     \caption{Salida del ejercicio 2}
%     \label{fig:ej2}
% \end{figure}

\newpage

% ============================================================
%  SECCIÓN 3 — Responsable: Integrante 2
%  Título: [Título del ejercicio según el taller]
% ============================================================

\section{[Título Sección 3]}

% ── Enunciado ───────────────────────────────────────────────
\subsection*{Enunciado}

\begin{tcolorbox}[colback=gray!8, colframe=gray!50, title=Ejercicio 3]
    Escriba aquí el enunciado del ejercicio 3.
\end{tcolorbox}

% ── Análisis / Planteamiento ────────────────────────────────
\subsection{Análisis del problema}

Descripción breve del enfoque o estrategia utilizada para resolver el ejercicio.

% ── Código ──────────────────────────────────────────────────
\subsection{Código}

\begin{lstlisting}[caption={Solución ejercicio 3}, label={lst:ej3}]
# Escribir aquí el código Python de la solución

def solucion():
    pass
\end{lstlisting}

% ── Respuesta / Resultado ───────────────────────────────────
\subsection{Resultados}

\begin{respuesta}
    Describir aquí los resultados obtenidos al ejecutar el programa.
\end{respuesta}

% ── (Opcional) Captura de pantalla ─────────────────────────
% \begin{figure}[h!]
%     \centering
%     \includegraphics[width=0.75\linewidth]{imagenes/resultado_ej3.png}
%     \caption{Salida del ejercicio 3}
%     \label{fig:ej3}
% \end{figure}

\newpage

% ============================================================
%  SECCIÓN 4 — Responsable: Integrante 2
%  Título: [Título del ejercicio según el taller]
% ============================================================

\section{[Título Sección 4]}

% ── Enunciado ───────────────────────────────────────────────
\subsection*{Enunciado}

\begin{tcolorbox}[colback=gray!8, colframe=gray!50, title=Ejercicio 4]
    Escriba aquí el enunciado del ejercicio 4.
\end{tcolorbox}

% ── Análisis / Planteamiento ────────────────────────────────
\subsection{Análisis del problema}

Descripción breve del enfoque o estrategia utilizada para resolver el ejercicio.

% ── Código ──────────────────────────────────────────────────
\subsection{Código}

\begin{lstlisting}[caption={Solución ejercicio 4}, label={lst:ej4}]
# Escribir aquí el código Python de la solución

def solucion():
    pass
\end{lstlisting}

% ── Respuesta / Resultado ───────────────────────────────────
\subsection{Resultados}

\begin{respuesta}
    Describir aquí los resultados obtenidos al ejecutar el programa.
\end{respuesta}

% ── (Opcional) Captura de pantalla ─────────────────────────
% \begin{figure}[h!]
%     \centering
%     \includegraphics[width=0.75\linewidth]{imagenes/resultado_ej4.png}
%     \caption{Salida del ejercicio 4}
%     \label{fig:ej4}
% \end{figure}

\newpage

% ============================================================
%  SECCIÓN 5 — Responsable: Integrante 3
%  Título: [Título del ejercicio según el taller]
% ============================================================

\section{[Título Sección 5]}

% ── Enunciado ───────────────────────────────────────────────
\subsection*{Enunciado}

\begin{tcolorbox}[colback=gray!8, colframe=gray!50, title=Ejercicio 5]
    Escriba aquí el enunciado del ejercicio 5.
\end{tcolorbox}

% ── Análisis / Planteamiento ────────────────────────────────
\subsection{Análisis del problema}

Descripción breve del enfoque o estrategia utilizada para resolver el ejercicio.

% ── Código ──────────────────────────────────────────────────
\subsection{Código}

\begin{lstlisting}[caption={Solución ejercicio 5}, label={lst:ej5}]
# Escribir aquí el código Python de la solución

def solucion():
    pass
\end{lstlisting}

% ── Respuesta / Resultado ───────────────────────────────────
\subsection{Resultados}

\begin{respuesta}
    Describir aquí los resultados obtenidos al ejecutar el programa.
\end{respuesta}

% ── (Opcional) Captura de pantalla ─────────────────────────
% \begin{figure}[h!]
%     \centering
%     \includegraphics[width=0.75\linewidth]{imagenes/resultado_ej5.png}
%     \caption{Salida del ejercicio 5}
%     \label{fig:ej5}
% \end{figure}

\newpage


% ── Conclusiones ────────────────────────────────────────────
% ============================================================
%  CONCLUSIONES — Responsable: Integrante 3
% ============================================================

\section{Conclusiones}

\subsection{Conclusiones generales}

Redactar aquí las conclusiones del taller. ¿Qué aprendieron? ¿Qué dificultades encontraron?
¿Cómo resolvieron los problemas planteados?

\subsection{Aportes individuales}

\begin{description}[leftmargin=2cm, style=nextline]
    \item[Integrante 1:] Describir brevemente el aporte y aprendizaje personal.
    \item[Integrante 2:] Describir brevemente el aporte y aprendizaje personal.
    \item[Integrante 3:] Describir brevemente el aporte y aprendizaje personal.
\end{description}

% ── Referencias (si aplica) ─────────────────────────────────
% \subsection*{Referencias}
% \begin{itemize}
%     \item Documentación oficial de Python: \url{https://docs.python.org/3/}
%     \item Recurso utilizado: \url{https://...}
% \end{itemize}


\end{document}
